\documentclass[12pt]{article}

\setlength{\parindent}{0pt}
\setcounter{secnumdepth}{3}
\usepackage[margin=1in]{geometry}

% packages
\usepackage{amsmath, amsfonts, amsthm, amssymb}
\usepackage{graphicx, float}
\usepackage{mathtools}
\usepackage{titlesec}
\usepackage{interval}
\usepackage{hyperref}
\usepackage{siunitx}
\usepackage{titling}
\usepackage{vwcol}
\usepackage{setspace}
\usepackage{empheq}
\usepackage{cancel}
\usepackage{esdiff}
\usepackage{multicol}
\usepackage{mdframed}
\usepackage{esdiff}
\usepackage{tikzsymbols}
\usepackage{multicol}
\usepackage{tikz}
\usepackage{varwidth}
\usepackage{pgfplots}
\usepackage{fancyhdr}
\usepackage{enumitem}
\usepackage{tcolorbox}

\intervalconfig {
	soft open fences
}

% header
\pagestyle{fancy}
\fancyhf{}
\lhead{Ethan Fan}
\rhead{PCH}
\cfoot{\thepage}

% section format
\titleformat{\section}
  {\bfseries}
  {}
  {0em}
  {}
\titlespacing*{\section}{0pt}{0pt}{0.3em}

\titleformat{\subsection}[runin]
  {\bfseries}
  {}
  {0em}
  {}
\titlespacing*{\subsection}{0pt}{0pt}{0em}

\titleformat{\subsubsection}[runin]
  {\itshape}
  {(\roman{subsubsection})}
  {0em}
  {}

% commands
\newcommand{\alignedintertext}[1]{%
  \noalign{%
    \vskip\belowdisplayshortskip
    \vtop{\hsize=\linewidth#1\par
    \expandafter}%
    \expandafter\prevdepth\the\prevdepth
  }%
}

\newcommand*{\paren}[1]{\ensuremath\left(#1\right)}
\newcommand*{\limit}[2][x]{\ensuremath{\displaystyle\lim_{#1\to#2}}}
\newcommand*{\Limit}[3][x]{\ensuremath{\displaystyle\lim_{#1\to#2}\left[#3\right]}}
\newcommand*{\deriv}[1][x]{\ensuremath{\dfrac{\mathrm{d}}{\mathrm{d}#1}}}
\newcommand*{\Deriv}[2][x]{\ensuremath{\dfrac{\mathrm{d}}{\mathrm{d}#1}\left[#2\right]}}
\newcommand{\dydx}{\frac{dy}{dx}}
\newcommand*{\iinteg}[2][x]{\ensuremath{\displaystyle\int #2\;\mathrm{d}#1}}
\newcommand*{\dinteg}[4][x]{\ensuremath{\displaystyle\int_{#2}^{#3}#4\;\mathrm{d}#1}}
\newcommand*{\abs}[1]{\ensuremath{\left|#1\right|}}
\newcommand*{\eps}{\varepsilon}
\newcommand*{\real}{\ensuremath\mathbb{R}}
\newcommand*{\integer}{\ensuremath\mathbb{Z}}
\newcommand*{\posint}{\ensuremath\mathbb{Z^+}}
\newcommand*{\cbrt}[1]{\ensuremath{\sqrt[3]{#1}}}
\newcommand*{\lheqzero}{\ensuremath{\underset{\text{L'H}}{\overset{\left[\frac00\right]}{=}}}}
\newcommand*{\lheqinfty}{\ensuremath{\underset{\text{L'H}}{\overset{\left[\frac{\infty}{\infty}\right]}{=}}}}
\newcommand{\tor}{\text{ or }}
\newcommand{\floor}[1]{\lfloor #1 \rfloor}

\newcommand{\vem}{\vspace{0.5em}}
\newcommand{\example}[1]{\section{Example #1}}
\newcommand{\problem}[1]{\section{Problem #1}}
\newcommand{\subproblem}[1]{\subsection{(#1)} \,}

\DeclareMathOperator{\dne}{DNE}
\DeclareMathOperator{\sgn}{sgn}

\newcommand{\definition}[1]{\begin{tcolorbox}[colback=red!5!white,colframe=red!75!black,parbox=false] #1 \end{tcolorbox}}

\newcommand{\theorem}[2]{\begin{tcolorbox}[title={#1},colback=blue!5!white,colframe=blue!75!black,parbox=false] #2 \end{tcolorbox}}

\allowdisplaybreaks
% \postdisplaypenalty=100000

% document
\begin{document}

\vspace*{0cm}

% opening
\begin{center}
    {\Large \bfseries Problem Set 65} \\[0.5cm]
    {\large Application of Derivatives 1} \\[0.35cm]
    {\normalsize April 13, 2026}
\end{center}

% problems

\problem{5}
Find the equations of all normal lines to the curve $y=x^3+2x+6$ that are parallel to the line $x+14y+4=0$.
\begin{align*}
    x+14y+4=0 & \implies \dydx =-\frac{1}{14}\\
    y=x^3+2x+6 & \implies \dydx = 3x^2+2\\
    3x^2+2=14 &\implies x^2=4 \implies x=\pm 2
\end{align*}
\[
\boxed{y-18=-\frac{1}{14}(x-2), \ y+6=-\frac{1}{14}(x+2)}
\]

\problem{6}
Consider the curve defined by the equation $x^3-y^2=0$. If the tangent line drawn at the point $P(4m^2,8m^2)$ is also a normal line to the same curve at some other point $Q$, find the value of $m$.
\begin{align*}
    y=x^\frac{3}{2} & \implies y'=\frac{3}{2}\sqrt{x}\\
    y-8m^3&=3m(x-4m^2)\\
    y=-x^\frac{3}{2} & \implies y'=-\frac{3}{2}\sqrt{x}\\
    -\frac{1}{3m}=-\frac{3}{2}\sqrt{x} &\implies x_2=\frac{4}{81m^2}\\
    y_2=-x_2^{\frac{3}{2}} &= -\frac{4}{81m^2}^{\frac{3}{2}} = - \frac{8}{729m^3} \\
    y-8m^3&=3m(x-4m^2)\\
    - \frac{8}{729m^3}-8m^3&=3m(\frac{4}{81m^2}-4m^2)\\
    - \frac{2}{729}-2m^6&=\frac{m^2}{27}-3m^6\\
    729m^6-27m^2-2&=0\\
    m^2 &= \frac{2}{9}\\
    & \hspace{-1.15em} \boxed{m=\pm \frac{\sqrt{2}}{3}}
\end{align*}

\problem{8}
Find the equations of all possible normal lines to the parabola $x^2=4y$ that pass through the point $(1,2)$.
\begin{align*}
    y=\frac{x^2}{4} \implies \dydx &=\frac{x}{2} \implies \frac{dx}{dy} = -\frac{2}{x}\\
    y-2&=-\frac{2}{x}(x-1) \\
    \frac{x^2}{4}-2&=-2+\frac{2}{x}\\
    x^3=8 & \implies x=2
\end{align*}
\[
\boxed{y-2=-(x-1)}
\]

\problem{9}
Find the coordinates of all points on the curve $y^2-2x^3-4y+8=0$ such that the tangent lines drawn at those points pass through the external point $(1,2)$.
\begin{gather*}
    2yy'-6x^2-4y'=0\\
    y'=\frac{3x^2}{y-2}\\
    (y-2)=\frac{3x^2}{y-2}(x-1)\\
    2x^3-4=3x^2(x-1)\\
    2x^3-4=3x^3-3x^2\\
    x^3-3x^2+4=0\\ 
    (x+1)(x-2)^2=0 \implies x=-1 \tor 2\\
    (y-2)^2=2x^3-4\\
    (y-2)^2=12\\
    y=2\pm 2\sqrt{3}\\
    \boxed{(2, 2+2\sqrt{3}),(2,2-2\sqrt{3})}
\end{gather*}

\problem{10}
Determine the condition under which the line $Ax+By=1$ is a normal line to the curve $a^{n-1}y=x^n$.
\begin{gather*}
    y=\frac{x^n}{a^{n-1}} \implies \dydx = n(\frac{x}{a})^{n-1}\implies \frac{dx}{dy}=\frac{1}{n} (\frac{a}{x})^{n-1}\\
    y=-\frac{A}{B}x+\frac{1}{B} \implies \frac{A}{B}=\frac{1}{n} (\frac{a}{x})^{n-1}\\
    Ax+B\frac{x^n}{a^{n-1}}=1 \implies x(A+\frac{B^2}{An})=1\\
    x=\frac{An}{A^2n+B^2}
\end{gather*}
Substituting $x$ and simplifying:
\begin{gather*}
    \frac{A}{B}=\frac{1}{n} (\frac{a}{x})^{n-1}\\
    \frac{A}{B}=\frac{1}{n} (\frac{a(A^2n+B^2)}{An})^{n-1}\\
    \boxed{(An)^n=Ba^{n-1}(A^2n+B^2)^{n-1}}
\end{gather*}

\problem{11}
Find the acute angle of intersection between the two curves $y = |x^2 - 1|$ and $y = |x^2 - 3|$.
\begin{gather*}
    |x^2-1|=|x^2-3|\\
    x^2-1=3-x^2\\
    x^2=2 \implies x=\sqrt{2}\\
    y_1'=(x^2-1)'=2x=2\sqrt{2}\\
    y_2'=(3-x^2)'=-2x=-2\sqrt{2}\\
    \tan \theta =|\frac{2\sqrt{2}-(-2\sqrt{2})}{1+(-8)}|=\frac{4\sqrt{2}}{7}\\
    \boxed{\theta = \arctan \paren{\frac{4\sqrt{2}}{7}}}
\end{gather*}

\problem{12}
Find the angle of intersection between the curves $2y^2 = x^3$ and $y^2 = 32x$ at every point where they
intersect.
\begin{gather*}
    y^2=\frac{x^3}{2}=32x\\
    x(x^2-64)=0\\
    x=0\implies y=0 \implies \theta =\frac{\pi}{2}\\   
    x=8 \implies y=16 \tor -16\\
    4yy_1'=3x^2 \implies y_1'=\frac{3x^2}{4y}\\
    2yy_2'=32 \implies y_2'=\frac{16}{y}\\
    \tan \theta = |\frac{3-1}{1+3\cdot 1}|=\frac{1}{2}\\
    \boxed{\theta = \arctan \paren{\frac{1}{2}} \tor \frac{\pi}{2}}
\end{gather*}

\problem{13}
Find the angle of intersection between the curves $f (x) = 2^x \ln x$ and $g(x) = x^{2x} - 1$ at their common point(s) of intersection.
\begin{gather*}
    2^x \ln x=x^{2x}-1 \implies x=1\\
    f'(x)=\ln 2 \cdot 2^x\cdot \ln x + 2^x \cdot \frac{1}{x} \implies f'(1)=2\\
    \ln y=2x\ln x\implies y'\cdot \frac{1}{y}=2\ln x+2\\
    g'(x)=x^{2x}(2 \ln x+2) \implies g'(1)=2\\
    \boxed{\theta = 0}
\end{gather*}

\problem{14}
Determine all possible values of the constant $a$ such that the ellipse $\frac{x^2}{a^2} + \frac{y^2}{4} = 1$ and the parabola $y^3 = 16x$ intersect orthogonally.
\begin{gather*}
    (\frac{x^2}{a^2} + \frac{y^2}{4})'=\frac{2x}
    {a^2}+\frac{2yy'}{4} \implies \dydx =-\frac{4x}{a^2y}\\
    3y^2y'=16 \implies \dydx =\frac{16}{3y^2}\\
    \frac{a^2y}{4x}=\frac{16}{3y^2} \\
    a^2=\frac{64x}{3y^3}\\
    a^2=\frac{4}{3}\\
    \boxed{a=\pm \frac{2}{\sqrt{3}}}
\end{gather*}

\problem{15}
A light ray originating from point $A(0, 2)$ strikes the curve $y = x^3 - x^2 + x - 1$ at the point $P (2, 5)$. Assuming the curve acts as a reflective surface, find the equation of the reflected ray.
\begin{gather*}
    \dydx =3x^2-2x+1 =12-4+1=9\\
    m_1=-\frac{1}{9}, m_2=\frac{3}{2}\\
    \tan \theta =|\frac{m_2-m_1}{1+m_1m_2}|=|\frac{\frac{27}{18}+\frac{2}{18}}{1-\frac{1}{6}}|=-\frac{29}{15}\\
    |\frac{m_3-m_1}{1+m_1m_3}|=\frac{29}{15}\\
    15(m_3+\frac{1}{9})=29(\frac{m_3}{9}-1)\\
    (15-\frac{29}{9})m_3=-29-\frac{5}{3}\\
    m_3=-\frac{276}{106}=-\frac{138}{53}\\
    \boxed{y-5=-\frac{138}{53}(x-2)}
\end{gather*}

\problem{16}
Find all possible real values of $p$ such that the equation $px^2 = \ln x$ has exactly one unique solution for $x$.
\begin{gather*}
    \frac{d}{dx}px^2=2px\\
    \frac{d}{dx}\ln x=\frac{1}{x}\\
    2px=\frac{1}{x} \implies p=\frac{1}{2x^2}\\
    px^2=\ln x \implies p=\frac{\ln x}{x^2}\\
    \ln x = \frac{1}{2}\implies x = \sqrt{e}\\
    \boxed{p=\frac{1}{2e}}
\end{gather*}

\problem{17}
Find the shortest distance between the line $y = x - 2$ and the parabola $y = x^2 + 3x + 2$.
\begin{gather*}
    \dydx = 1\\
    \dydx =2x+3=1 \implies x=-1\\
    y=x+1\\
    \boxed{d=\frac{3}{\sqrt{2}}}
\end{gather*}

\problem{18}
Find the minimum value of $(x_1 -x_2)^2 + \paren{\dfrac{x_1^2}{20} - \sqrt{(17 - x_2)(x_2 - 13)}}^2$, where $x_1 \in  \real^+ $ and $x_2 \in (13, 17)$.
\begin{gather*}
    y=\frac{x^2}{20}\\
    y=\sqrt{4-(x-15)^2} \implies y^2+(x-15)^2=2^2\\
    f(x)=(x-15)^2+(\frac{x^2}{20}-0)^2\\
    f'(x)=2(x-15)+\frac{x^3}{100}\\
    x^3+200x-3000=0 \implies x=10\\
    D=\sqrt{(10-15)^2+5^2}-2=5\sqrt{2}-2\\
    \boxed{D^2=54-20\sqrt{2}}
\end{gather*}

\problem{19}
Consider the curve defined by the equation:
\[
y^2=4a\paren{x+a\sin \frac{x}{a}}
\]
Prove that all points on this curve where the tangent is parallel to the $x$-axis must lie on the parabola
$y^2 = 4ax$.
\begin{gather*}
    2yy'=4a(1+a\cos \frac{x}{a} \cdot \frac{1}{a})\\
    y'=\frac{2a}{y}(1+\cos \frac{x}{a})=0\\
    \cos \frac{x}{a}=-1 \implies \sin  \frac{x}{a}=0\\
    y^2=4a(x+a\cdot 0)\\
    \boxed{y^2=4ax}
\end{gather*}

\end{document}

